\section{Discussion and Limitations}
\label{sec:discussion}

\subsection{What the Completed Evidence Establishes}
The results answer the central research question at the level of completed configurations. Frozen representations support metric aerial depth when paired with trainable task heads. Dense decoding and camera-conditioned readout produce lower reported depth error than the corresponding small probes. The final combination of frozen DINOv3 context features, a frozen native multiview geometry provider, and VoxDet decoding yields a repeated occupied-space result, while its semantic macro score remains low. This separation is the principal empirical finding: recovering occupied structure does not by itself solve fine-grained aerial recognition.

Architecture attribution is more limited. The multiview treatment differs from the historical FoundationSSC control in decoding, objective, and optimization. The two monocular backbone records differ in pretrained depth quality and realized validation budget. No factorial experiment separately varies all these components with repeated matched controls. We consequently describe consistent directions and complete-treatment differences, avoiding claims that VoxNT, camera metadata, or MVS alone explains a specific percentage-point gain.

\subsection{Geometry Quality and Semantic Recognition}
A scalar depth metric measures visible surface error. SSC consumes geometry through projection, feature accumulation, volumetric fusion, and completion, then evaluates an extended voxel space. Errors near boundaries can place features in the wrong voxel even when average depth error is modest. Conversely, a posterior that spreads mass across plausible depths may support occupied coverage despite an imperfect expectation. These differences help explain why the depth and semantic rankings are not identical, but the present data do not measure their causal contributions.

The final method explicitly retains a categorical geometry posterior. This provides a richer interface than a single expected depth, especially when matching evidence admits multiple candidates. However, categorical probabilities should not be interpreted as calibrated epistemic or obstacle-risk uncertainty. The network and adapter assign normalized mass, while reliable risk estimation would require additional validation of calibration, coverage, and behavior under distribution shift. Out-of-range exclusion and invalid-ray masking preserve the adapter's geometric contract; they do not establish the accuracy of the retained posterior.

Semantic recognition has a separate supervision bottleneck. A class-balanced objective reduces domination by empty and large occupied regions, but it cannot create visual evidence or annotation support for every category. Zero true positives for person, bicycle, and cable persist despite explicit occupied-class weighting and eligible directional supervision. The failure could involve classification, voxel placement, object visibility, or their interaction. Further diagnosis requires per-target predictions and confusion analysis tied to the evaluated manifest rather than another aggregate score alone.

\subsection{Reading Long-Tail Scores Correctly}
The standard 21-class macro average uses a fixed denominator and assigns zero contribution to classes absent from the evaluated split. It is a stable reporting rule, not proof that recognition of every category was tested. Per-class N/A labels identify absent ground-truth support, while an observed zero remains a failure on a represented category. The same convention must carry through captions and comparisons.

The support figure describes the test distribution. Without split-wide train/validation histograms, we cannot establish whether a rare test category was also rare or absent in training. Appearance ambiguity, visibility, voxel size, geometry, and annotation can influence IoU alongside frequency. Per-target predictions and complete split counts are needed to distinguish these mechanisms.

\subsection{Repetition, Data Identity, and Generalization}
Aligning model and sampler seeds makes each independent run reproducible under its training contract. The three-seed sample standard deviation reports variability across those runs, with all other stated settings held fixed. It is not a confidence interval over environments and does not turn a one-run historical control into a repeated baseline. A paired comparison should repeat the control with the same declared seeds and compare differences on the same evaluation targets.

All experiments concern one aerial dataset with two test scenes and flight heights represented in training. The paper does not establish generalization to other camera systems, weather, terrain, altitudes, or geographic regions. Nearby frames also share content, so a large frame or voxel count should not be confused with a large number of independent environments. Additional datasets and scene-level analyses are needed before drawing broader conclusions about aerial foundation-model transfer.

\subsection{Implications for Robotic Use}
The completed system provides a useful occupied-space research result, but its deployment implications must follow the measured tasks. A landing or inspection planner may need correct recognition of people, vehicles, cables, and traversable surfaces. Persistent failures in represented categories prevent claiming such semantic reliability from the headline completion score. A dense completed volume can contribute hypotheses to a planner, while newly observed evidence should update those hypotheses; the current experiments do not evaluate a closed-loop planner or certify an unobserved region as free.

The five-view geometry input creates additional practical requirements: synchronized or consistently indexed observations, calibrated poses, source-view selection, and sufficient image overlap. Pose error, changing objects, weak texture, and occlusion can affect correspondence. Their sensitivity is not measured in this study. A benchmark configuration with retained calibration is not a demonstration that the same geometry will remain accurate with online UAV pose estimates.

Compute is another unresolved dimension. Freezing large encoders reduces training adaptation but leaves their inference cost. The multiview provider processes higher-resolution images, and SSC constructs a dense voxel representation. Offline training on eight H100s does not quantify latency, memory, energy, or throughput on embedded hardware. The published efficiency of an upstream decoder cannot be transferred to this hybrid without measuring the complete frontend, geometry, fusion, and decoding path.

\subsection{Focused Next Experiments}
The most informative follow-up is a repeated matched comparison of the geometry-family control and final decoder treatment, with one target manifest and a shared validation search window. This would separate variability of the complete treatment from variability of the baseline and permit paired descriptive differences. A subsequent component study could vary full-resolution offset supervision, the occupancy/semantic objective, and depth-gradient adaptation while preserving the remaining contract.

For depth, primary result recovery and shared-target evaluation should precede additional architectural claims. Equal-budget repeated image heads would clarify decoder and metadata effects, while a common context-valid target set would permit a temporal comparison. For semantics, per-target confusion analysis and complete split histograms would help distinguish sparse supervision from geometric or appearance failures. These are proposed experiments, not unfinished numeric results inserted into the current tables.

\section{Conclusion}
\label{sec:conclusion}
We studied frozen image/video foundation representations for aerial metric depth and semantic scene completion on OccuFly. The completed depth comparison supports trainable dense readout and explicit camera conditioning, with provenance limits stated for retained MDE records. Combining frozen DINOv3 features and native calibrated multiview geometry with VoxDet-style voxel decoding produces a stable occupied-space result across three independent runs. Standard 21-class semantic evaluation and class-wise support show that this geometric progress coexists with substantial unresolved category recognition. The evidence supports the complete transfer configuration and motivates matched repeated controls, stronger per-target provenance, and rare-class diagnosis before deployment claims are made.
